\subsection{连通集在连续映射下的像}

命题 4. 设 A 是拓扑空间 X 的连通子集，f 是 X 到拓扑空间 X' 中的一个连续映射。则 \(f(\mathbf{A})\) 是连通的。

设 \(f(\mathbf{A})\) 不连通。则存在两个集合 \(\mathbf{M}'\), \(\mathbf{N}'\)，它们在 \(f(\mathbf{A})\) 中是开的，且构成 \(f(\mathbf{A})\) 的一个分拆；于是 \(\mathbf{A} \cap \overline{f}^{1}(\mathbf{M}')\) 与 \(\mathbf{A} \cap \overline{f}^{1}(\mathbf{N}')\) 在 \(\mathbf{A}\) 中是开的，且构成 \(\mathbf{A}\) 的一个分拆；这与 \(\mathbf{A}\) 连通的假设相矛盾。

连通集在连续映射下的逆像未必连通；例如考虑一个离散空间到一个由单独一点组成的空间中的映射。

由命题 4 我们得到不连通空间的另一个刻画：

命题 5. 拓扑空间 X 不连通的必要充分条件是：存在一个 X 到含有多于一个点的离散空间上的连续满射。

由命题 4，该条件是充分的。反之，若 X 不连通，则存在两个非空不相交的开子集 A, B，其并为 X，而 X 到两个元素的离散空间 \(\{a, b\}\) 上的映射 f（由 \(f(A) = \{a\}\) 与 \(f(B) = \{b\}\) 定义）是连续的。

